%% APPENDIX C -- the velocity-space line mask and its cost.
%% Owner: r10-frame.  Carries \label{app:linemask} (cited by
%% 04_method.tex) and \label{sec:linecost}.
\section{The velocity-space line mask, and what it costs}
\label{app:linemask}

Attribution decides the candidate list and the mask decides attribution, so
the mask's definition and its price are both given here.

\emph{Definition.} A one-channel coincidence test is too tight for
circumstellar gas: a native channel spans a fraction of a km\,s$^{-1}$ while
real emission is shifted by many channels. The tolerance follows from the
kinematic budget of the gas rather than from the instrument --- Keplerian
motion where these molecules survive ($\lesssim$13\,km\,s$^{-1}$ beyond
10\,au) and linewidths $\lesssim$1\,km\,s$^{-1}$ --- and is
$\pm\MaskHalfKms$\,km\,s$^{-1}$ around each catalogued transition.

\emph{The frame, which is the star's own.} Searched frequencies are
topocentric as the correlator delivers them. Before any transition is
compared with a crossing, that crossing's frequency is carried to the
barycentre using its own block's pointing and mid-time and then to the
stellar rest frame using the star's systemic velocity. The barycentric term
runs from \FrBaryLo{} to \FrBaryHi\,km\,s$^{-1}$ across the survey, so a
tolerance applied in the observed frame would spend
\FrBarySwing\,km\,s$^{-1}$ of a $\pm\MaskHalfKms$\,km\,s$^{-1}$ budget on
the Earth's motion, and the offsets of individual crossings differ between
the two frames by up to \MaskFrameShiftKms\,km\,s$^{-1}$. The
\FrEpochNEpoch{} \FrEpochStar{} epochs show the transform working in the
only way that can be checked: \FrEpochTopoKms\,km\,s$^{-1}$ apart in the
observed frame, they agree to \FrEpochStelKms\,km\,s$^{-1}$ in the star's.

The mask holds \FrNTrans{} transitions of \FrNSpec{} species at full
JPL/CDMS precision \citep{Pickett1998,CDMS2005}. The species criterion is
the one objective rule statable in advance: the transitions ALMA's band
definitions are built around --- the CO isotopologues, HCN, HCO$^{+}$, CS,
SiO, CN, H$_2$CO, SO and [C\,\textsc{i}] --- which dominate the emission of
disc-selected fields. Sulphur monoxide belongs there for the same reason
carbon monosulphide does: \FrNSOWord{} of its transitions fall in the
intervals searched, it is a routine component of the Band~6 and Band~7
setups these fields were observed with, and sulphur bearing molecules are
expected wherever the carbon bearing ones already masked are. It is not a
formality --- the strongest crossing that passes the spatial screen,
\FrCpStar{} at $T_\star=\VnCpT$, lies \FrCpSODv\,km\,s$^{-1}$ from
\FrSOName{} in its own star's frame (\S\ref{sec:chaingap}).

\emph{The half-width is not a tuned parameter, and the counts rather than
the argument say so.} Table~\ref{tab:maskrobust} carries the whole
disposition chain at half-widths from $\pm\FrLadderLo$ to
$\pm\FrLadderHi$\,km\,s$^{-1}$. One count moves, and it is worth naming:
narrowing to $\pm\FrLadderLo$\,km\,s$^{-1}$ releases \FrNarrowMoveStar{}
from attribution at \FrNarrowMoveDv\,km\,s$^{-1}$ from
\FrNarrowMoveLine{} --- a crossing that does not outrank its own control
ring either way. Everything the result turns on is flat: \FrNRankPass{}
crossing outranks every control at every half-width in that range, and it is
the same crossing each time, while the fraction of searched bandwidth the
mask removes rises from \FrMaskPctLo{} to \FrMaskPctHi{} per cent. Widening
the tube therefore buys nothing and costs coverage, and narrowing it
releases an identified carbon monoxide line of this survey's own.

\begin{table}
\centering
\footnotesize
\caption{The disposition chain against the mask half-width, in the stellar
rest frame. All \NCross{} crossings are classified at every half-width: the
second column is what the mask costs in searched bandwidth and the rest is
what it decides, ending in how many unattributed crossings outrank all
\NCtrl{} controls and how many of \emph{those} recur at a later epoch.
Adopted half-width in bold.}
\label{tab:maskrobust}
%% GENERATED by maskframe_v411.py -- do not hand-edit.
\begin{tabular}{@{}lrrrrr@{}}
\hline
Half-width & band & line- & unatt- & outrank & recur- \\
(km\,s$^{-1}$) & masked & attributed & ributed & all & ring \\
 & (per cent) & & & controls & \\
\hline
$\pm20$ & 2.1 & 18 & 38 & 1 & 0 \\
$\pm25$ & 2.6 & 19 & 37 & 1 & 0 \\
\textbf{$\pm50$} & \textbf{4.8} & \textbf{19} & \textbf{37} & \textbf{1} & \textbf{0} \\
$\pm100$ & 8.7 & 19 & 37 & 1 & 0 \\
\hline
\end{tabular}

\end{table}

\emph{The species list is a choice, and widening it is not free.} The
alternative selection --- every Splatalogue transition over the
\UnionIntervals{} searched islands with $E_{\rm u}\le\MaskDbEuK$\,K and
$\log A_{ij}\ge\MaskDbLogA$, which is \MaskDbTrans{} transitions of
\MaskDbCarriers{} species --- masks \MaskDbPctUnion{} per cent of the union
against \FrMaskPct{} per cent here, and it does not leave the dispositions
alone. \FrFullNFlip{} of the \NCross{} crossings change side under it,
\FrFullNFlipRankWord{} of them a crossing that outranks every control:
\FrFullFlipRankStar{} falls \FrFullFlipRankDv\,km\,s$^{-1}$ from a
vibrationally excited \FrFullFlipRankChem{} transition, which nobody would
predict toward a quiet main-sequence G star. That is the calibration and not
a defect in either list: at these frequencies a uniformly random frequency
in one of these windows lies within $\pm\MaskHalfKms$\,km\,s$^{-1}$ of
something in the wide harvest most of the time. ``Coincident with a known
molecular line'' is a statement about a catalogue rather than a property of
a frequency, and a survey quoting one should quote the occupancy of the list
it used.

\subsection{The cost of the mask}
\label{sec:linecost}

Masking removes \FrMaskLostGHz\,GHz, \FrMaskPct{} per cent of the
\FrMaskUnionGHz\,GHz unique-frequency union, CN, CO and SO dominating. It
is not uniform: the loss reaches \FrMaskWorstPct{} per cent in
Band~\FrMaskWorstBand{} and stays below 2\,per cent in Bands~3--5, and it
is tabulated band by band in the data release. The fine-channel experiment
pays the larger fractional price, its coverage being concentrated in the
line-tuned windows.

What masking buys is classification rather than detection: a carrier at a
line frequency is as detectable as anywhere else but cannot be told from the
line, because the astrophysical prior at a catalogued transition is
overwhelming. The consequences for the limits are in
\S\ref{sec:excludes}, and a design that searches every frequency and
attaches a \emph{lower} prior to a line-coincident event
\citep{CocconiMorrison1959} is in \S\ref{sec:future}; the released
channels support that second search.
